\section{团队组织与成员分工}\label{sec:division}
\subsection{小组之间的职责}
第一组负责本次撰写的人形机器人与方式二部分。其余采用“第$i$组”代称；下表为方便协作，暂将第二、三、四组分别对应方式一、三、四，需由对应小组确认后定稿。仓库部分历史 README 使用“本组负责方式三”的叙述，不能据此推断第一组成员身份。
\begin{longtable}{p{28mm}p{66mm}p{43mm}}
\caption{全队功能分工（人员信息待核实）}\\\toprule
功能组 & 工作范围与交付 & 协同接口/负责人\\\midrule\endfirsthead
\toprule 功能组 & 工作范围与交付 & 协同接口/负责人\\\midrule\endhead
第一组：人形机器人 & 动作录制与组合、抓取/放置、相机调试、视觉闭环路线、运行日志 & 目标背部标记、任务状态；姓名待填\\
第二组：方式一（暂定） & AR直接跟随实际代码、标签尺寸、速度控制、失去标记处置；核对与课程固定航路点方式的区别 & \code{follow1}；姓名待填\\
第一组：方式二 & 地图采集、AMCL、AR转地图、轨迹目标、move\_base与速度门控 & \code{follow2}；姓名待填\\
第三组：方式三（暂定） & RGB-D目标检测、身份关联、实时避障与独立验收 & \code{follow3}；姓名待填\\
第四组：方式四（暂定） & 前一台TurtleBot位姿/轨迹获取、队列间距与末车停靠 & \code{follow4}；姓名待填\\
联调与报告 & 网络/TF核对、各方式独立性确认、数据汇总、视频与报告合并 & 总负责人、各组对接人待填\\\bottomrule
\end{longtable}

\subsection{本小组内部职责登记}
目前未提供成员姓名及个人贡献，以下是可认领的工作项，\textbf{不是对实际个人贡献的认定}。一人可承担多项，一项可由多人共同完成；贡献比例仅在成员协商后填写。
\begin{longtable}{p{23mm}p{26mm}p{48mm}p{37mm}}
\caption{本小组成员与贡献登记}\\\toprule
姓名/学号 & 职责方向 & 具体工作及交付证据 & 贡献说明/确认\\\midrule\endfirsthead
\toprule 姓名/学号 & 职责方向 & 具体工作及交付证据 & 贡献说明/确认\\\midrule\endhead
待填写 & 动作与机体 & 录制抓取/放置，组合抱物步态，检查姿态衔接；对应动作文件与提交号 & 待成员核实\\
待填写 & 视觉与状态机 & 相机维护、标记检测、路线切换、窗口与日志；对应代码与运行记录 & 待成员核实\\
待填写 & 建图与导航 & Kinect/扫描调试、建图、AMCL、代价地图与DWA参数；地图及配置 & 待成员核实\\
待填写 & 跟随与协同 & AR地图目标、轨迹回退、速度门控、网络及任务状态；联调日志 & 待成员核实\\
待填写 & 测试与报告 & 数据采集、统计、图片归档、章节编写与排版；原始记录及文稿 & 待成员核实\\\bottomrule
\end{longtable}

\subsection{其他小组内部职责登记规范}
方式一、三、四的章节各预留一张成员分工表。各组填写“姓名、学号、具体任务、交付证据、合作边界”，避免只写“参与调试”或“负责报告”。最终由各组负责人核对，并在全队分工表登记对接人。
